% HiddenPower — итерация 5, накопительная версия V06.
% Глава: «Стохастические дифференциальные уравнения и диффузионные модели».
% Источники: S. Särkkä, A. Solin; Y. Song et al.

\chapter{Стохастические дифференциальные уравнения и диффузионные модели}
\label{chap:stochastic-differential-equations}

В главе~\ref{chap:dynamics-differentiable-models} состояние $x(t)$ определялось начальным условием и законом динамики. Такая модель описывает одну траекторию при фиксированных параметрах. В реальной конструкции часть воздействий и свойств известна лишь приближённо: нагрузка меняется, материал неоднороден, геометрия имеет допуски, а вычислительная модель не учитывает все масштабы процесса. Тогда одного значения $x(t)$ недостаточно. Нужно описывать множество возможных состояний и вероятности их появления.

Основная цепочка этой главы имеет вид
\[
\begin{aligned}
\text{неопределённость}
&\longrightarrow
\text{случайная траектория},\\
\text{случайная траектория}
&\longrightarrow
\text{эволюция распределения},\\
\text{распределение}
&\longrightarrow
\text{обратная динамика}
\longrightarrow
\text{генерация}.
\end{aligned}
\]
Сначала вводятся случайные процессы и броуновское движение. Затем шум включается в дифференциальное уравнение через интеграл Ито, а распределение состояния описывается уравнением Фоккера--Планка. В последней части главы эта конструкция приводит к score-функции, обратному SDE и диффузионной генерации.

\section{Случайные процессы и броуновское движение}
\label{sec:stochastic-processes-brownian-motion}

\subsection{От одной траектории к семейству возможных траекторий}

Пусть $\Omega$ обозначает множество возможных реализаций неопределённости. Элемент $\omega\in\Omega$ фиксирует один конкретный сценарий: одну последовательность возмущений, одну реализацию производственных отклонений или один случайный выбор начального состояния.

\begin{definition}
Случайным процессом со значениями в $\R^d$ называется семейство случайных векторов
\begin{equation}
\label{eq:stochastic-process-definition-v01}
X=\{X_t\}_{t\in I},
\qquad
X_t:\Omega\to\R^d,
\end{equation}
индексированное временем $t$ из интервала $I$.
\end{definition}

При фиксированном $t$ величина $X_t$ является случайным вектором. При фиксированном $\omega$ получается функция времени
\begin{equation}
\label{eq:sample-path-v01}
t\longmapsto X_t(\omega),
\end{equation}
которая называется реализацией или выборочной траекторией процесса. Это различие принципиально: процесс задаёт распределение над траекториями, а одна численно построенная кривая показывает только один исход.

Детерминированная модель из главы~\ref{chap:dynamics-differentiable-models} является частным случаем. Если для всех $\omega$ выполняется $X_t(\omega)=x(t)$, то случайность отсутствует и все реализации совпадают. В стохастической модели одинаковое начальное описание может приводить к разным траекториям.

\begin{application}[состояние конструкции как случайный процесс]
Пусть
\begin{equation}
\label{eq:stochastic-design-state-v01}
X_t=
\begin{pmatrix}
G_t\\
M_t\\
H_t\\
R_t
\end{pmatrix}
\end{equation}
содержит параметры геометрии $G_t$, материала $M_t$, толщины $H_t$ и пористости или параметров решётки $R_t$. Индекс $t$ может обозначать физическое время, уровень приложенной нагрузки или искусственное время зашумления. Одна реализация $X_t(\omega)$ описывает один возможный сценарий эволюции конструкции, а распределение $X_t$ характеризует неопределённость по всему ансамблю сценариев.
\end{application}

Само по себе множество траекторий ещё не даёт вычислимой модели. Нужно уметь описывать распределения процесса в выбранные моменты времени.

\subsection{Распределения, среднее и ковариация}

Для каждого набора времён $t_1,\ldots,t_k$ случайный процесс задаёт совместное распределение вектора
\begin{equation}
\label{eq:finite-dimensional-distribution-v01}
\begin{pmatrix}
X_{t_1}\\
\vdots\\
X_{t_k}
\end{pmatrix}
\in\R^{kd}.
\end{equation}
Эти распределения называют конечномерными. Они отвечают на вопросы о значениях процесса в нескольких моментах и содержат информацию о временной зависимости.

Первое числовое описание процесса даёт среднее
\begin{equation}
\label{eq:process-mean-v01}
m(t)=\mathbbE[X_t]\in\R^d.
\end{equation}
Оно показывает центр распределения состояния, но не определяет разброс. Для двух моментов $s$ и $t$ вводится ковариационная матрица
\begin{equation}
\label{eq:process-covariance-v01}
C(s,t)
=
\mathbbE\!\left[
\bigl(X_s-m(s)\bigr)
\bigl(X_t-m(t)\bigr)^{\trans}
\right].
\end{equation}
При $s=t$ получаем матрицу ковариации состояния
\begin{equation}
\label{eq:state-covariance-v01}
P(t)=C(t,t).
\end{equation}
Она симметрична и неотрицательно определена, поскольку для любого $v\in\R^d$
\[
v^{\trans}P(t)v
=
\mathbbE\!\left[
\bigl(v^{\trans}(X_t-m(t))\bigr)^2
\right]
\geq0.
\]

Связь с линейной алгеброй непосредственна. Если
\[
P(t)u_i=\lambda_i u_i,
\qquad
\lambda_i\geq0,
\]
то собственный вектор $u_i$ задаёт главное направление разброса, а $\sqrt{\lambda_i}$ --- характерный масштаб отклонения вдоль него. В координатах параметров конструкции ковариационный эллипсоид показывает не только отдельные допуски, но и совместные направления неопределённости.

Нулевая ковариация не всегда означает независимость. Для совместно гауссовских величин это верно, но в общем случае зависимость может быть нелинейной и не отражаться во втором моменте. Поэтому среднее и ковариация полностью описывают процесс только при дополнительных предположениях.

\begin{definition}
Процесс $X_t$ называется гауссовским, если для любых $t_1,\ldots,t_k$ вектор~\eqref{eq:finite-dimensional-distribution-v01} имеет многомерное нормальное распределение.
\end{definition}

Для гауссовского процесса функции $m(t)$ и $C(s,t)$ определяют все конечномерные распределения. Это делает гауссовские процессы особенно удобными: бесконечномерный объект задаётся средним и ковариационным ядром, а вычисления на конечной сетке сводятся к операциям с векторами и матрицами.

\begin{connection}[связь с предыдущими главами]
Матрица $P(t)$ использует спектральный аппарат главы~\ref{chap:linear-algebra}. Производные функций от состояния и преобразование плотностей опираются на главу~\ref{chap:multivariable-calculus}. Если допустимые конструкции образуют многообразие, то координатный шум должен согласовываться с геометрией главы~\ref{chap:manifolds-differential-geometry}. Наконец, временная эволюция случайного состояния продолжает детерминированную динамику главы~\ref{chap:dynamics-differentiable-models}.
\end{connection}

\subsection{Броуновское движение}

Базовым источником непрерывного шума служит броуновское движение, или процесс Винера. В прикладном изложении Särkkä и Solin оно задаётся через непрерывность, независимые гауссовские приращения и начальное значение.

\begin{definition}
Стандартным $r$-мерным броуновским движением называется процесс $W_t\in\R^r$, $t\geq0$, для которого выполнены условия
\begin{equation}
\label{eq:brownian-origin-v01}
W_0=0,
\end{equation}
траектории $t\mapsto W_t(\omega)$ непрерывны, а для любых $0\leq s<t$
\begin{equation}
\label{eq:brownian-increment-v01}
W_t-W_s
\sim
\mathcal N\!\left(0,(t-s)I_r\right).
\end{equation}
Приращения на непересекающихся временных интервалах независимы.
\end{definition}

Последнее условие означает, что новые случайные изменения не содержат информации о прошлых приращениях. Оно не означает независимость самих значений $W_s$ и $W_t$: более позднее значение включает более раннее.

Из определения сразу следуют
\begin{equation}
\label{eq:brownian-mean-variance-v01}
\mathbbE[W_t]=0,
\qquad
\mathbbE[W_tW_t^{\trans}]=tI_r.
\end{equation}
Типичный масштаб перемещения растёт как $\sqrt{t}$, а не как $t$. Действительно, стандартное отклонение каждой координаты равно $\sqrt{t}$.

Для $0\leq s\leq t$ запишем
\[
W_t=W_s+(W_t-W_s).
\]
Приращение $W_t-W_s$ независимо от $W_s$ и имеет нулевое среднее. Поэтому
\begin{align}
\mathbbE[W_sW_t^{\trans}]
&=
\mathbbE\!\left[
W_sW_s^{\trans}
\right]
+
\mathbbE\!\left[
W_s(W_t-W_s)^{\trans}
\right]
\\
&=
sI_r+0.
\end{align}
С учётом симметрии по $s$ и $t$ получаем ковариационную функцию
\begin{equation}
\label{eq:brownian-covariance-v01}
C_W(s,t)=\min(s,t)I_r.
\end{equation}
Эта формула показывает, какая часть случайной истории является общей для двух моментов времени.

Более общий процесс с диффузионной матрицей $Q\succeq0$ имеет приращения
\begin{equation}
\label{eq:brownian-general-covariance-v01}
B_t-B_s
\sim
\mathcal N\!\left(0,(t-s)Q\right),
\qquad
C_B(s,t)=\min(s,t)Q.
\end{equation}
Если $Q=LL^{\trans}$, то можно положить $B_t=LW_t$. Фактор $L$ преобразует независимые стандартные направления шума в коррелированные возмущения физических параметров.

\begin{computationalnote}[моделирование приращений]
На сетке $0=t_0<t_1<\cdots<t_N$ броуновскую траекторию строят рекурсией
\begin{equation}
\label{eq:brownian-discrete-simulation-v01}
B_{t_{k+1}}
=
B_{t_k}
+
\sqrt{\Delta t_k}\,Lz_k,
\qquad
z_k\sim\mathcal N(0,I_r),
\end{equation}
где $\Delta t_k=t_{k+1}-t_k$, а векторы $z_k$ независимы. Тогда
\[
\operatorname{Cov}
\bigl(B_{t_{k+1}}-B_{t_k}\bigr)
=
\Delta t_kLL^{\trans}
=
\Delta t_kQ.
\]
Множитель $\sqrt{\Delta t_k}$ обязателен. Замена его на $\Delta t_k$ дала бы дисперсию порядка $\Delta t_k^2$ и неправильный предел при измельчении сетки.
\end{computationalnote}

\subsection{Законченный расчёт вероятности приращения}

Пусть отклонение толщины в упрощённой модели описывается
\begin{equation}
\label{eq:thickness-brownian-model-v01}
H_t=H_0+\sigma W_t,
\qquad
\sigma=0.2\ \text{мм}/\sqrt{\text{с}},
\end{equation}
где $W_t$ --- стандартное одномерное броуновское движение. Требуется найти вероятность того, что за интервал длиной $9$ секунд изменение толщины по модулю не превысит $0.5$ мм.

Дано
\[
H_{t+9}-H_t
=
\sigma(W_{t+9}-W_t).
\]
По формуле~\eqref{eq:brownian-increment-v01}
\[
W_{t+9}-W_t\sim\mathcal N(0,9),
\]
следовательно,
\begin{equation}
\label{eq:thickness-increment-distribution-v01}
H_{t+9}-H_t
\sim
\mathcal N\!\left(0,\sigma^2\cdot9\right)
=
\mathcal N(0,0.36\ \text{мм}^2).
\end{equation}
Стандартное отклонение приращения равно $0.6$ мм. После стандартизации
\begin{align}
\mathbbP\!\left(
\abs{H_{t+9}-H_t}\leq0.5\ \text{мм}
\right)
&=
\mathbbP\!\left(
\abs{Z}\leq\frac{0.5}{0.6}
\right)
\\
&=
2\Phi(0.8333)-1
\\
&\approx0.595,
\end{align}
где $Z\sim\mathcal N(0,1)$, а $\Phi$ --- функция распределения стандартной нормальной величины. Вероятность составляет около $59.5\%$.

Расчёт показывает порядок действий: сначала определяется распределение приращения, затем вычисляется его дисперсия на нужном интервале, после чего событие переводится к стандартной нормальной величине. Одной величины $\sigma$ недостаточно без длительности интервала: разброс растёт как $\sigma\sqrt{\Delta t}$.

\subsection{Белый шум и предел классической производной}

Броуновские траектории непрерывны, но почти наверное нигде не дифференцируемы. На малом интервале длиной $\Delta t$ приращение имеет порядок $\sqrt{\Delta t}$, поэтому формальное отношение
\[
\frac{W_{t+\Delta t}-W_t}{\Delta t}
\]
имеет типичный масштаб $1/\sqrt{\Delta t}$ и не стабилизируется при $\Delta t\to0$.

Белый шум записывают как формальную производную
\begin{equation}
\label{eq:white-noise-formal-derivative-v01}
\xi(t)=\frac{\dd W_t}{\dd t},
\end{equation}
но это не обычная функция времени. Запись полезна для физической интуиции, однако выражение $\xi(t)\,\dd t$ нужно понимать через приращение $\dd W_t$. Именно поэтому классический интеграл Римана недостаточен и в следующем разделе будет введён интеграл Ито.

\begin{caution}[что нельзя делать с белым шумом]
Нельзя подставлять белый шум в ОДУ как обычную ограниченную функцию и затем применять стандартные правила дифференцирования. Такая запись скрывает выбор стохастического интеграла. Для линейных моделей эвристика иногда приводит к верному результату, но в нелинейном случае разные способы предельного перехода могут давать разные уравнения.
\end{caution}

\subsection{Прикладной смысл и ограничения модели}

В генеративном проектировании броуновское движение удобно не потому, что производственные ошибки буквально являются движением частицы. Оно даёт управляемую модель непрерывного накопления множества малых независимых возмущений. Матрица $Q$ задаёт направления и интенсивность зашумления, а время регулирует общий масштаб неопределённости.

Если $X_0$ является случайной конструкцией из набора данных, то процесс
\begin{equation}
\label{eq:additive-forward-diffusion-preview-v01}
X_t=X_0+B_t
\end{equation}
постепенно размывает исходное распределение. В работе Song и соавторов эта идея обобщается до прямого SDE, которое переводит распределение данных в простое шумовое распределение. До раздела V06 достаточно различать отдельную зашумлённую конструкцию и распределение ансамбля.

Модель имеет жёсткие ограничения. Аддитивный гауссовский шум не сохраняет положительность толщины, интервал допустимой пористости и нелинейные геометрические ограничения. Простое отсечение значений меняет распределение и создаёт искусственные массы на границе. Если допустимые состояния лежат на многообразии, одинаковый евклидов шум в разных координатах может иметь разный геометрический смысл. Масштабы координат также критичны: один и тот же коэффициент шума нельзя без нормировки применять одновременно к длине в миллиметрах и безразмерной доле материала.

Поэтому перед использованием процесса нужно явно зафиксировать четыре объекта: кодировку состояния, единицы измерения, ковариацию $Q$ и область допустимых значений. Только после этого случайная траектория получает физический смысл. В следующих версиях дрейф будет описывать направленную динамику, диффузия --- интенсивность случайных возмущений, а эволюция плотности свяжет отдельные траектории с распределением конструкций.

\section{Интеграл Ито и стохастическое дифференциальное уравнение}
\label{sec:ito-integral-sde}

\subsection{Почему обычного интеграла недостаточно}

Формальная запись белого шума из~\eqref{eq:white-noise-formal-derivative-v01} подсказывает уравнение
\begin{equation}
\label{eq:white-noise-driven-formal-v02}
\frac{\dd X_t}{\dd t}
=
f(X_t,t)+G(X_t,t)\,\xi(t).
\end{equation}
Однако правая часть не является обычной функцией времени: броуновская траектория почти наверное не имеет производной. Поэтому~\eqref{eq:white-noise-driven-formal-v02} служит только мнемонической записью. Математический смысл задаётся интегральным уравнением.

Обычный вклад динамики интегрируется по времени:
\begin{equation}
\label{eq:drift-integral-v02}
\int_{t_0}^{t} f(X_s,s)\,\dd s.
\end{equation}
Для шумового вклада требуется интеграл относительно броуновского движения
\begin{equation}
\label{eq:stochastic-integral-required-v02}
\int_{t_0}^{t} G(X_s,s)\,\dd W_s.
\end{equation}
Разница между ними не формальна. На интервале длиной $\Delta t$ приращение времени равно $\Delta t$, а приращение броуновского движения имеет типичный размер $\sqrt{\Delta t}$. Поэтому шумовой вклад нельзя получить заменой случайной функции в обычном интеграле Римана.

Пусть задано разбиение
\[
t_0=\tau_0<\tau_1<\cdots<\tau_n=t.
\]
Для обычного интеграла выбор точки внутри малого интервала не меняет предел. При интегрировании по броуновской траектории выбор точки существенен. Интеграл Ито фиксирует левый конец интервала; тем самым коэффициент шума вычисляется до появления нового приращения.

\begin{caution}[не смешивать белый шум и броуновское движение]
Белый шум $\xi(t)$ не является случайной величиной с конечной дисперсией в каждой точке времени. Вычислимым объектом служит приращение $\Delta W_k=W_{\tau_{k+1}}-W_{\tau_k}$. Выражения вида $\xi(t)^2$ или $h(\xi(t))$ для нелинейной функции $h$ нельзя использовать как обычные функции без отдельного предельного определения.
\end{caution}

\subsection{Определение интеграла Ито}

Сначала рассмотрим одномерное броуновское движение и ступенчатый процесс
\begin{equation}
\label{eq:simple-adapted-process-v02}
H_s
=
\sum_{k=0}^{n-1}
H_k\,\mathbf 1_{(\tau_k,\tau_{k+1}]}(s).
\end{equation}
Коэффициент $H_k$ может зависеть от случайной истории до момента $\tau_k$, но не от будущего приращения $W_{\tau_{k+1}}-W_{\tau_k}$. Такое условие называют адаптированностью. В прикладной модели оно запрещает выбирать текущий коэффициент шума после того, как уже стало известно будущее возмущение.

\begin{definition}
Интеграл Ито ступенчатого адаптированного процесса~\eqref{eq:simple-adapted-process-v02} определяется равенством
\begin{equation}
\label{eq:ito-integral-simple-v02}
\int_{t_0}^{t} H_s\,\dd W_s
:=
\sum_{k=0}^{n-1}
H_k
\bigl(W_{\tau_{k+1}}-W_{\tau_k}\bigr).
\end{equation}
Для общего процесса $H_s$ интеграл определяется как предел таких сумм в среднем квадратичном, если
\begin{equation}
\label{eq:ito-square-integrability-v02}
\mathbb E\!\left[
\int_{t_0}^{t} H_s^2\,\dd s
\right]<\infty.
\end{equation}
\end{definition}

Сходимость в среднем квадратичном означает
\[
\mathbb E\!
\left[
\abs{I_n-I}^2
\right]
\longrightarrow0.
\]
Это сильнее сходимости средних значений и контролирует типичный квадрат ошибки. Полная теория требует вероятностного пространства и фильтрации, но для дальнейших вычислений достаточно трёх свойств: коэффициент не использует будущее, выполнено условие~\eqref{eq:ito-square-integrability-v02}, а интеграл строится левыми суммами.

Для векторного процесса $W_t\in\R^r$ и матричного коэффициента $G_s\in\R^{d\times r}$ интеграл
\begin{equation}
\label{eq:ito-integral-matrix-v02}
\int_{t_0}^{t}G_s\,\dd W_s
\in\R^d
\end{equation}
понимается покоординатно. Условие квадратичной интегрируемости принимает вид
\begin{equation}
\label{eq:ito-matrix-integrability-v02}
\mathbb E\!\left[
\int_{t_0}^{t}\norm{G_s}_{\mathrm F}^2\,\dd s
\right]<\infty,
\end{equation}
где $\norm{G}_{\mathrm F}^2$ равно сумме квадратов элементов матрицы.

\subsection{Среднее, ковариация и изометрия Ито}

Пусть сначала $H_s$ имеет вид~\eqref{eq:simple-adapted-process-v02}. Каждое новое приращение броуновского движения имеет нулевое среднее и независимо от информации, доступной в левом конце интервала. Поэтому
\begin{equation}
\label{eq:ito-integral-zero-mean-v02}
\mathbb E\!\left[
\int_{t_0}^{t}H_s\,\dd W_s
\right]=0.
\end{equation}

Рассмотрим квадрат суммы~\eqref{eq:ito-integral-simple-v02}. Смешанные слагаемые с разными индексами имеют нулевое среднее, а для диагональных слагаемых
\[
\mathbb E\!\left[
H_k^2(\Delta W_k)^2
\right]
=
\mathbb E[H_k^2]\,\Delta\tau_k.
\]
Переходя к пределу, получаем изометрию Ито
\begin{equation}
\label{eq:ito-isometry-v02}
\mathbb E\!\left[
\left(
\int_{t_0}^{t}H_s\,\dd W_s
\right)^2
\right]
=
\mathbb E\!\left[
\int_{t_0}^{t}H_s^2\,\dd s
\right].
\end{equation}
Она переводит случайный интеграл в обычный интеграл от квадрата коэффициента. Именно это равенство обеспечивает устойчивость определения в среднем квадратичном.

Если $G(s)$ является детерминированной матрицей, то случайный вектор
\begin{equation}
\label{eq:deterministic-matrix-ito-v02}
Y_t=\int_{t_0}^{t}G(s)\,\dd W_s
\end{equation}
имеет нулевое среднее и ковариацию
\begin{equation}
\label{eq:ito-integral-covariance-v02}
\operatorname{Cov}(Y_t)
=
\int_{t_0}^{t}G(s)G(s)^{\trans}\,\dd s.
\end{equation}
Матрица $G(s)G(s)^{\trans}$ задаёт скорость накопления ковариации. Это продолжает факторизацию $Q=LL^{\trans}$ из V01, но теперь интенсивность и направления шума могут меняться со временем.

\begin{connection}[связь с линейной алгеброй]
Матрица $G$ отображает $r$ независимых шумовых координат в $d$ координат состояния. Ранг $G$ показывает число непосредственно возбуждаемых направлений, а собственные векторы $GG^{\trans}$ задают главные направления мгновенного разброса. Если $\operatorname{rank}G<d$, шум локально действует только в подпространстве, хотя дрейф может переносить его влияние на остальные координаты.
\end{connection}

\subsection{SDE как интегральное уравнение}

\begin{definition}
Стохастическим дифференциальным уравнением Ито называется запись
\begin{equation}
\label{eq:general-ito-sde-v02}
\dd X_t
=
f(X_t,t)\,\dd t
+
G(X_t,t)\,\dd W_t,
\qquad
X_{t_0}=X_0,
\end{equation}
которая означает интегральное равенство
\begin{equation}
\label{eq:general-ito-integral-equation-v02}
X_t
=
X_0
+
\int_{t_0}^{t}f(X_s,s)\,\dd s
+
\int_{t_0}^{t}G(X_s,s)\,\dd W_s.
\end{equation}
Функция $f:\R^d\times I\to\R^d$ называется дрейфом, а $G:\R^d\times I\to\R^{d\times r}$ --- коэффициентом диффузии.
\end{definition}

Дрейф описывает направленное изменение среднего масштаба порядка $\Delta t$. Диффузия преобразует случайное приращение порядка $\sqrt{\Delta t}$. Условное приращение на малом интервале имеет приближённые характеристики
\begin{align}
\mathbb E[\Delta X_t\mid X_t=x]
&\approx
f(x,t)\,\Delta t,
\label{eq:local-drift-moment-v02}
\\
\operatorname{Cov}(\Delta X_t\mid X_t=x)
&\approx
G(x,t)G(x,t)^{\trans}\,\Delta t.
\label{eq:local-diffusion-moment-v02}
\end{align}
Поэтому матрицу
\begin{equation}
\label{eq:diffusion-tensor-v02}
D(x,t)=G(x,t)G(x,t)^{\trans}
\end{equation}
можно понимать как локальную скорость роста ковариации.

Размерности коэффициентов различаются. Если $X$ измеряется в единицах состояния, то $f$ имеет единицы состояния на единицу времени, а $G$ --- состояния на квадратный корень из времени. Эта проверка часто обнаруживает ошибочно выбранный масштаб шума.

В работе Song и соавторов прямой процесс зашумления записан в той же форме, причём коэффициент диффузии зависит от времени. Позднее обучаемая score-функция изменит дрейф обратной динамики, но само прямое SDE уже полностью задаётся функциями $f$ и $G$ и не требует нейронной сети.

\begin{caution}[дифференциальная запись является сокращением]
Нельзя алгебраически сокращать $\dd t$, считать $\dd W_t/\dd t$ обычной функцией или применять классическое правило цепочки к~\eqref{eq:general-ito-sde-v02}.

Все преобразования должны сохранять интегральный смысл SDE. Исправленное правило цепочки будет выведено в V03 как формула Ито.
\end{caution}

\subsection{Существование и единственность решения}

Решение~\eqref{eq:general-ito-sde-v02} называется сильным, если при заданных начальном состоянии и броуновском движении процесс $X_t$ строится на той же случайной истории. Потраекторная единственность означает: два решения с одинаковыми $X_0$ и $W_t$ совпадают почти наверное. Именно такой смысл нужен при численном моделировании, когда одна и та же последовательность случайных приращений должна определять одну траекторию.

Достаточные условия имеют ту же структуру, что теорема существования для ОДУ из главы~04. Пусть $K>0$. Тогда для всех $x,y\in\R^d$ и $t\in[t_0,T]$ требуется
\begin{equation}
\label{eq:sde-global-lipschitz-v02}
\norm{f(x,t)-f(y,t)}
+
\norm{G(x,t)-G(y,t)}_{\mathrm F}
\leq
K\norm{x-y},
\end{equation}
а также выполнена линейная оценка роста
\begin{equation}
\label{eq:sde-linear-growth-v02}
\norm{f(x,t)}^2
+
\norm{G(x,t)}_{\mathrm F}^2
\leq
K^2\bigl(1+\norm{x}^2\bigr).
\end{equation}
Тогда на конечном интервале существует единственное сильное решение, и его второй момент конечен.

Условия~\eqref{eq:sde-global-lipschitz-v02}--\eqref{eq:sde-linear-growth-v02} являются достаточными, но не необходимыми. Их нарушение не доказывает отсутствие решения. Однако без отдельного анализа возможны несколько решений, взрыв траектории за конечное время или потеря конечности моментов.

Для нейросетевого дрейфа или диффузии гладкость архитектуры сама по себе не даёт глобальной оценки. Большие веса могут сделать константу Липшица огромной, а полиномиальный рост --- нарушить~\eqref{eq:sde-linear-growth-v02}. На практике проверяют диапазон состояний, ограничивают спектральные нормы слоёв, контролируют величину коэффициентов и отслеживают траектории, покидающие рабочую область.

\subsection{Законченный расчёт для SDE с постоянными коэффициентами}

Рассмотрим скалярное уравнение
\begin{equation}
\label{eq:additive-sde-v02}
\dd X_t
=
a\,\dd t
+
\sigma\,\dd W_t,
\qquad
X_0=x_0,
\end{equation}
где $a$ и $\sigma$ постоянны. Подставляем коэффициенты в интегральную форму~\eqref{eq:general-ito-integral-equation-v02}:
\begin{align}
X_t
&=
x_0
+
\int_0^t a\,\dd s
+
\int_0^t \sigma\,\dd W_s
\\
&=
x_0+at+\sigma W_t.
\label{eq:additive-sde-solution-v02}
\end{align}
Поскольку $W_t\sim\mathcal N(0,t)$,
\begin{equation}
\label{eq:additive-sde-distribution-v02}
X_t
\sim
\mathcal N\!\left(x_0+at,\sigma^2t\right).
\end{equation}
Следовательно,
\begin{equation}
\label{eq:additive-sde-moments-v02}
\mathbb E[X_t]=x_0+at,
\qquad
\operatorname{Var}(X_t)=\sigma^2t.
\end{equation}
Дрейф сдвигает центр распределения линейно по времени, а диффузия увеличивает стандартное отклонение как $\sigma\sqrt t$.

Пусть $X_t$ обозначает толщину элемента в миллиметрах,
\[
x_0=2,
\qquad
a=-0.03\ \text{мм}/\text{с},
\qquad
\sigma=0.1\ \text{мм}/\sqrt{\text{с}}.
\]
Требуется найти распределение через $25$ секунд и вероятность нефизического события $X_{25}<0$.

По~\eqref{eq:additive-sde-distribution-v02}
\begin{align}
\mathbb E[X_{25}]
&=
2-0.03\cdot25
=
1.25\ \text{мм},
\\
\operatorname{Var}(X_{25})
&=
0.1^2\cdot25
=
0.25\ \text{мм}^2.
\end{align}
Стандартное отклонение равно $0.5$ мм, поэтому
\begin{align}
\mathbb P(X_{25}<0)
&=
\mathbb P\!\left(
Z<\frac{0-1.25}{0.5}
\right)
\\
&=
\Phi(-2.5)
\approx
0.0062.
\end{align}
Вероятность отрицательной толщины составляет около $0.62\%$. Она мала, но не равна нулю; значит, даже простая модель с правдоподобными средним и дисперсией нарушает физическое ограничение.

Для шага длиной $\Delta t$ из~\eqref{eq:additive-sde-solution-v02} точно следует
\begin{equation}
\label{eq:additive-sde-exact-step-v02}
X_{t+\Delta t}
=
X_t
+
a\Delta t
+
\sigma\sqrt{\Delta t}\,Z,
\qquad
Z\sim\mathcal N(0,1).
\end{equation}
В данном линейном случае это точный переход, а не численная аппроксимация. Для переменных коэффициентов та же структура станет методом Эйлера--Маруямы в V05.

\begin{application}[что нужно проверить перед использованием SDE конструкции]
Координаты состояния должны быть нормированы и иметь однозначные единицы. Дрейф обязан отражать направленный физический или искусственный процесс, а матрица $D=GG^{\trans}$ --- допустимые направления и скорость накопления неопределённости. Если толщина, плотность или пористость ограничены, аддитивное SDE в исходных координатах обычно непригодно: требуется преобразование переменной, отражающая граница либо динамика в параметризации, где ограничения выполняются автоматически. Наконец, стохастический шум моделирует выбранный вид вариабельности; он сам по себе не отделяет шум данных от неопределённости модели.
\end{application}

\section{Формула Ито и уравнение Фоккера--Планка}
\label{sec:ito-formula-fokker-planck}

\subsection{Почему обычное правило цепочки перестаёт работать}

Пусть $X_t$ удовлетворяет скалярному SDE
\begin{equation}
\label{eq:scalar-sde-v03}
\dd X_t
=
f(X_t,t)\,\dd t
+
g(X_t,t)\,\dd W_t.
\end{equation}
Для гладкой функции $\varphi(x,t)$ требуется найти динамику величины $Y_t=\varphi(X_t,t)$. В детерминированном случае достаточно первого порядка формулы Тейлора. Для SDE это неверно, поскольку броуновское приращение имеет масштаб
\[
\Delta W_t=O(\sqrt{\Delta t}).
\]
Следовательно,
\[
(\Delta W_t)^2=O(\Delta t),
\]
и квадратичный член формулы Тейлора имеет тот же порядок, что и дрейф. Отбрасывать его нельзя.

Ключевой факт выражается через квадратичную вариацию. Для разбиения $0=t_0<t_1<\cdots<t_N=t$ рассмотрим сумму
\begin{equation}
\label{eq:quadratic-variation-sum-v03}
Q_N(t)
=
\sum_{k=0}^{N-1}
\bigl(W_{t_{k+1}}-W_{t_k}\bigr)^2.
\end{equation}
При измельчении разбиения эта сумма сходится к $t$. В символической записи используют правила
\begin{equation}
\label{eq:ito-multiplication-table-v03}
(\dd W_t)^2=\dd t,
\qquad
\dd W_t\,\dd t=0,
\qquad
(\dd t)^2=0.
\end{equation}
Это не алгебра бесконечно малых чисел, а краткая запись предельных соотношений для стохастических сумм. Правила~\eqref{eq:ito-multiplication-table-v03} допустимы внутри вычислений Ито, но не заменяют интегральное определение из раздела~\ref{sec:ito-integral-sde}.

Простейший пример показывает появление поправки. Для $Y_t=W_t^2$ имеем
\[
\dd Y_t
=
2W_t\,\dd W_t
+
(\dd W_t)^2
=
2W_t\,\dd W_t+\dd t.
\]
Поэтому обычная цепочка $\dd(W_t^2)=2W_t\,\dd W_t$ теряет детерминированный член. Именно эта потеря приводит к неверным средним, дисперсиям и плотностям.

\subsection{Одномерная формула Ито}

\begin{theorem}[формула Ито в одном измерении]
Пусть $X_t$ решает~\eqref{eq:scalar-sde-v03}, а функция $\varphi:\R\times[0,T]\to\R$ имеет непрерывные производные $\partial_t\varphi$, $\partial_x\varphi$ и $\partial_{xx}\varphi$. Тогда процесс $Y_t=\varphi(X_t,t)$ удовлетворяет
\begin{equation}
\label{eq:ito-formula-scalar-v03}
\begin{aligned}
\dd \varphi(X_t,t)
={}&
\left(
\partial_t\varphi
+
f\,\partial_x\varphi
+
\frac12 g^2\,\partial_{xx}\varphi
\right)(X_t,t)\,\dd t
\\
&+
\left(g\,\partial_x\varphi\right)(X_t,t)\,\dd W_t.
\end{aligned}
\end{equation}
\end{theorem}

Формула получается из разложения второго порядка
\[
\dd\varphi
\approx
\partial_t\varphi\,\dd t
+
\partial_x\varphi\,\dd X_t
+
\frac12\partial_{xx}\varphi\,(\dd X_t)^2.
\]
Подставляя~\eqref{eq:scalar-sde-v03} и используя~\eqref{eq:ito-multiplication-table-v03}, находим
\[
(\dd X_t)^2
=
\bigl(f\,\dd t+g\,\dd W_t\bigr)^2
=
g^2\,\dd t.
\]
Отсюда появляется член $\tfrac12g^2\partial_{xx}\varphi$. Он зависит от кривизны функции $\varphi$ и интенсивности шума. Если $g=0$, формула Ито переходит в обычное правило цепочки.

\begin{connection}[связь с многомерным анализом]
Производная $\partial_x\varphi$ измеряет локальную чувствительность преобразования, а $\partial_{xx}\varphi$ --- его кривизну. В детерминированной динамике первый порядок управляет изменением функции вдоль траектории. В стохастической динамике гессиан также влияет на средний дрейф, потому что случайные отклонения исследуют окрестность точки сразу в обоих направлениях.
\end{connection}

Формула Ито нужна не только для замены переменной. Она позволяет строить уравнения для моментов. Если взять $\varphi(x)=x$, то получаем исходное SDE. Если взять $\varphi(x)=x^2$, то
\begin{equation}
\label{eq:ito-square-general-v03}
\dd(X_t^2)
=
\bigl(2X_tf(X_t,t)+g^2(X_t,t)\bigr)\,\dd t
+
2X_tg(X_t,t)\,\dd W_t.
\end{equation}
После усреднения стохастический интеграл исчезает при стандартных условиях интегрируемости, и возникает уравнение для второго момента.

\subsection{Многомерная формула и генератор SDE}

Пусть $X_t\in\R^d$ удовлетворяет
\begin{equation}
\label{eq:multidimensional-sde-v03}
\dd X_t
=
f(X_t,t)\,\dd t
+
G(X_t,t)\,\dd W_t,
\end{equation}
где $W_t\in\R^r$. Обозначим локальную диффузионную матрицу
\begin{equation}
\label{eq:diffusion-matrix-v03}
D(x,t)=G(x,t)G(x,t)^{\trans}\in\R^{d\times d}.
\end{equation}
Для гладкой скалярной функции $\varphi:\R^d\times[0,T]\to\R$ формула Ито принимает вид
\begin{equation}
\label{eq:ito-formula-vector-v03}
\begin{aligned}
\dd\varphi(X_t,t)
={}&
\left[
\partial_t\varphi
+
\nabla\varphi^{\trans}f
+
\frac12\operatorname{tr}
\left(D\,\nabla^2\varphi\right)
\right](X_t,t)\,\dd t
\\
&+
\left[\nabla\varphi^{\trans}G\right](X_t,t)\,\dd W_t.
\end{aligned}
\end{equation}
Здесь $\nabla^2\varphi$ --- гессиан по координатам состояния. След матричного произведения разворачивается как
\[
\operatorname{tr}(D\nabla^2\varphi)
=
\sum_{i,j=1}^{d}
D_{ij}\,
\frac{\partial^2\varphi}{\partial x_i\partial x_j}.
\]
Диагональные элементы $D$ отвечают за дисперсию отдельных координат, а внедиагональные --- за совместные шумовые направления.

\begin{definition}
Инфинитезимальным генератором SDE~\eqref{eq:multidimensional-sde-v03} называется дифференциальный оператор
\begin{equation}
\label{eq:generator-v03}
(\mathcal A_t\varphi)(x)
=
\nabla\varphi(x)^{\trans}f(x,t)
+
\frac12
\operatorname{tr}
\left(D(x,t)\nabla^2\varphi(x)\right).
\end{equation}
Для функции, зависящей от времени, к правой части добавляется $\partial_t\varphi$.
\end{definition}

Генератор описывает мгновенную скорость изменения среднего значения тестовой функции. Формально
\begin{equation}
\label{eq:generator-expectation-v03}
\frac{\dd}{\dd t}\mathbb E[\varphi(X_t)]
=
\mathbb E[(\mathcal A_t\varphi)(X_t)].
\end{equation}
Он действует на наблюдаемую $\varphi$, тогда как уравнение Фоккера--Планка будет действовать на плотность распределения. Эти две точки зрения связаны сопряжённостью операторов.

\subsection{От SDE к эволюции плотности}

Пусть распределение $X_t$ имеет плотность $p(x,t)$. Тогда для тестовой функции $\varphi$ выполняется
\begin{equation}
\label{eq:expectation-density-v03}
\mathbb E[\varphi(X_t)]
=
\int_{\R^d}\varphi(x)p(x,t)\,\dd x.
\end{equation}
Дифференцируя по времени и используя~\eqref{eq:generator-expectation-v03}, получаем слабое равенство
\begin{equation}
\label{eq:fpk-weak-form-v03}
\int_{\R^d}
\varphi(x)\,\partial_t p(x,t)\,\dd x
=
\int_{\R^d}
(\mathcal A_t\varphi)(x)p(x,t)\,\dd x.
\end{equation}
Перенос производных с $\varphi$ на $p$ выполняется интегрированием по частям. Если плотность и поток достаточно быстро убывают на бесконечности либо заданы согласованные граничные условия, краевые члены исчезают.

\begin{theorem}[уравнение Фоккера--Планка]
Пусть коэффициенты SDE~\eqref{eq:multidimensional-sde-v03} достаточно регулярны, а распределение решения имеет гладкую плотность $p(x,t)$. Тогда
\begin{equation}
\label{eq:fokker-planck-v03}
\partial_t p(x,t)
=
-\sum_{i=1}^{d}
\frac{\partial}{\partial x_i}
\bigl(f_i(x,t)p(x,t)\bigr)
+
\frac12
\sum_{i,j=1}^{d}
\frac{\partial^2}{\partial x_i\partial x_j}
\bigl(D_{ij}(x,t)p(x,t)\bigr).
\end{equation}
Начальное условие задаётся плотностью $p(x,t_0)=p_0(x)$.
\end{theorem}

Первый член переносит вероятность полем дрейфа. Второй распределяет её по направлениям, заданным $D$. При постоянной изотропной диффузии $D=\sigma^2I_d$ уравнение упрощается до
\begin{equation}
\label{eq:fpk-isotropic-v03}
\partial_t p
=
-\nabla\cdot(fp)
+
\frac{\sigma^2}{2}\Delta p.
\end{equation}
Здесь дивергенция описывает перенос плотности, а лапласиан --- её сглаживание. Это прямое продолжение градиента, дивергенции и оператора Лапласа из глав~\ref{chap:multivariable-calculus} и~\ref{chap:manifolds-differential-geometry}.

Уравнение~\eqref{eq:fokker-planck-v03} сохраняет полную вероятность при нулевом потоке через границу:
\[
\frac{\dd}{\dd t}
\int p(x,t)\,\dd x
=0.
\]
На ограниченной области это свойство зависит от граничного условия. Поглощающая граница удаляет массу, отражающая сохраняет её внутри области, а периодическая связывает противоположные части границы. Поэтому постановка FPK без указания области и границы неполна.

\subsection{Законченный расчёт: броуновское движение с дрейфом}

Рассмотрим SDE с постоянными коэффициентами
\begin{equation}
\label{eq:drifted-brownian-sde-v03}
\dd X_t
=
a\,\dd t
+
\sigma\,\dd W_t,
\qquad
X_0=x_0.
\end{equation}
В V02 было получено точное решение
\[
X_t=x_0+at+\sigma W_t.
\]
Теперь выведем ту же эволюцию на уровне плотности. Здесь
\[
f(x,t)=a,
\qquad
D(x,t)=\sigma^2.
\]
Подстановка в~\eqref{eq:fokker-planck-v03} даёт
\begin{equation}
\label{eq:advection-diffusion-v03}
\partial_t p
=
-a\,\partial_x p
+
\frac{\sigma^2}{2}\partial_{xx}p,
\qquad
p(x,0)=\delta(x-x_0).
\end{equation}
Решение имеет нормальную плотность
\begin{equation}
\label{eq:drifted-brownian-density-v03}
p(x,t)
=
\frac{1}{\sqrt{2\pi\sigma^2t}}
\exp\left(
-\frac{(x-x_0-at)^2}{2\sigma^2t}
\right).
\end{equation}
Проверим параметры на конкретных данных:
\[
x_0=2,
\qquad
a=-0.03,
\qquad\sigma=0.1,
\qquad t=25.
\]
Центр равен
\[
\mu_{25}=x_0+a t=2-0.03\cdot25=1.25,
\]
а дисперсия
\[
v_{25}=\sigma^2t=0.1^2\cdot25=0.25.
\]
Следовательно,
\begin{equation}
\label{eq:worked-density-v03}
p(x,25)
=
\frac{1}{\sqrt{0.5\pi}}
\exp\left(
-\frac{(x-1.25)^2}{0.5}
\right).
\end{equation}
При дрейфе центр плотности переместился от $2$ к $1.25$, а диффузия увеличила стандартное отклонение до $0.5$. Интегральное решение SDE описывает отдельные траектории, а FPK описывает тот же процесс как перенос и расширение всей плотности.

Для проверки среднего умножим~\eqref{eq:advection-diffusion-v03} на $x$ и проинтегрируем. После интегрирования по частям получаем
\[
\frac{\dd}{\dd t}\mathbb E[X_t]=a.
\]
При начальном условии $\mathbb E[X_0]=x_0$ отсюда следует $\mathbb E[X_t]=x_0+at$. Аналогично для второго момента формула Ито~\eqref{eq:ito-square-general-v03} даёт
\[
\frac{\dd}{\dd t}\operatorname{Var}(X_t)=\sigma^2,
\]
то есть $\operatorname{Var}(X_t)=\sigma^2t$. Потраекторное и распределительное описания согласуются.

\subsection{Практический смысл и ограничения}

В генеративном проектировании прямое SDE задаёт контролируемое разрушение структуры распределения конструкций. Если $X_0$ кодирует геометрию, материал, толщину и параметры решётки, то $p(x,t)$ показывает, как распределяется ансамбль этих состояний после уровня шума $t$. Дрейф может сжимать или масштабировать пространство состояний, а диффузия сглаживает плотность и постепенно стирает локальные особенности исходных данных.

Уравнение Фоккера--Планка важно концептуально: оно связывает локальное правило движения каждой траектории с глобальным изменением распределения. Именно плотность $p(x,t)$ появится в V06 внутри score-функции
\[
\nabla_x\log p(x,t).
\]
Пока эта функция не вычисляется и не аппроксимируется сетью; здесь устанавливается только источник её временной зависимости.

Прямое решение FPK практически ограничено размерностью. После дискретизации каждой координаты число узлов растёт экспоненциально, поэтому сеточные методы быстро становятся непригодными для векторов геометрии и полей высокой размерности. В таких задачах используют выборочные траектории SDE, аналитические переходные плотности специальных моделей или обучаемые приближения локальных характеристик распределения.

Существование сильного решения SDE ещё не гарантирует существование гладкой плотности. Вырожденная матрица $D$ может распределять шум только по подпространству; ограничения могут создавать границы и сингулярности; дискретные или категориальные параметры не имеют обычной плотности относительно объёма в $\R^d$. Если допустимые состояния лежат на многообразии, евклидово SDE может немедленно вывести траекторию наружу. Тогда нужны координатно согласованная диффузия, проекция или отдельная геометрическая модель.

\begin{application}[контроль постановки распределительной модели]
Перед использованием FPK нужно определить пространство состояния, единицы и масштаб каждой координаты, начальную плотность, область, граничные условия и матрицу $D=GG^{\trans}$. Затем проверяют сохранение полной вероятности, положительность плотности и согласование её моментов с симуляциями SDE. Без этих проверок гладкая численная картинка может описывать не исходную стохастическую модель, а артефакт сетки или неверно заданной границы.
\end{application}

\section{Линейные SDE, процесс Орнштейна--Уленбека, моменты и ковариации}
\label{sec:linear-sde-ou-moments}

Нелинейное SDE обычно не допускает замкнутого решения ни для траектории, ни для плотности. Линейный случай устроен проще: состояние является линейным преобразованием начального условия и броуновского движения. Поэтому при гауссовском начальном распределении процесс остаётся гауссовским, а его закон полностью задаётся средним и ковариацией.

Линейные модели важны не только как учебные примеры. Они описывают малые отклонения около рабочего режима, дают локальное приближение нелинейной динамики и позволяют проверить численный код по точным формулам. Процесс Орнштейна--Уленбека показывает основной механизм: случайные возмущения расширяют распределение, а устойчивый дрейф возвращает состояние к опорному уровню.

\subsection{Линейное SDE и матрица перехода}

Рассмотрим процесс $X_t\in\R^d$, заданный уравнением
\begin{equation}
\label{eq:linear-sde-v04}
\dd X_t
=
\bigl(F(t)X_t+u(t)\bigr)\,\dd t
+
G(t)\,\dd W_t,
\qquad
X_{t_0}=X_0.
\end{equation}
Здесь $F(t)\in\R^{d\times d}$ задаёт линейную динамику, $u(t)\in\R^d$ --- детерминированное воздействие, $G(t)\in\R^{d\times r}$ --- направления и интенсивности шума, а $W_t$ --- стандартное $r$-мерное броуновское движение.

Матрица перехода $\Phi(t,s)$ определяется задачей
\begin{equation}
\label{eq:transition-matrix-v04}
\frac{\partial}{\partial t}\Phi(t,s)
=
F(t)\Phi(t,s),
\qquad
\Phi(s,s)=I_d.
\end{equation}
Она переносит состояние однородной детерминированной системы из момента $s$ в момент $t$. Для переменной матрицы $F(t)$ в общем случае нельзя заменить $\Phi(t,s)$ обычной экспонентой от интеграла: матрицы $F(t_1)$ и $F(t_2)$ могут не коммутировать.

Интегральное решение~\eqref{eq:linear-sde-v04} имеет вид
\begin{equation}
\label{eq:linear-sde-solution-v04}
X_t
=
\Phi(t,t_0)X_0
+
\int_{t_0}^{t}\Phi(t,\tau)u(\tau)\,\dd\tau
+
\int_{t_0}^{t}\Phi(t,\tau)G(\tau)\,\dd W_\tau.
\end{equation}
Первое слагаемое переносит начальную неопределённость, второе описывает известное воздействие, третье накапливает случайные возмущения. Формула является стохастическим аналогом вариации постоянных из главы~\ref{chap:dynamics-differentiable-models}.

\begin{connection}[связь с линейной алгеброй и динамикой]
Матрица $\Phi(t,s)$ действует на векторы состояния так же, как фундаментальная матрица линейной ODE. Её сингулярные числа показывают усиление конечных возмущений, а собственные значения постоянной матрицы $F$ определяют устойчивость отдельных мод. Шум не заменяет эту динамику: он непрерывно возбуждает те направления, которые затем усиливаются или подавляются оператором перехода.
\end{connection}

\subsection{Уравнения для среднего и ковариации}

Обозначим
\begin{equation}
\label{eq:mean-covariance-v04}
m(t)=\mathbb E[X_t],
\qquad
P(t)=\mathbb E\!\left[(X_t-m(t))(X_t-m(t))^{\trans}\right].
\end{equation}
Пусть $X_0$ независимо от будущих приращений $W_t-W_{t_0}$ и имеет конечный второй момент. Усреднение~\eqref{eq:linear-sde-v04} даёт
\begin{equation}
\label{eq:linear-mean-ode-v04}
\dot m(t)=F(t)m(t)+u(t),
\qquad
m(t_0)=m_0.
\end{equation}
Стохастический интеграл не входит в среднее, поскольку его ожидание равно нулю.

Для ковариации формула Ито приводит к матричному уравнению
\begin{equation}
\label{eq:linear-covariance-ode-v04}
\dot P(t)
=
F(t)P(t)+P(t)F(t)^{\trans}+G(t)G(t)^{\trans},
\qquad
P(t_0)=P_0.
\end{equation}
Первые два члена переносят уже существующий разброс линейной динамикой. Последний член добавляет новую ковариацию от шума. Решения равны
\begin{align}
\label{eq:linear-mean-solution-v04}
m(t)
&=
\Phi(t,t_0)m_0
+
\int_{t_0}^{t}\Phi(t,\tau)u(\tau)\,\dd\tau,
\\
\label{eq:linear-covariance-solution-v04}
P(t)
&=
\Phi(t,t_0)P_0\Phi(t,t_0)^{\trans}
+
\int_{t_0}^{t}
\Phi(t,\tau)G(\tau)G(\tau)^{\trans}
\Phi(t,\tau)^{\trans}\,\dd\tau.
\end{align}
Обе части~\eqref{eq:linear-covariance-solution-v04} неотрицательно определены. Поэтому точная формула сохраняет допустимую матричную структуру ковариации.

Если $X_0$ имеет гауссовское распределение, то~\eqref{eq:linear-sde-solution-v04} является суммой гауссовских величин и
\begin{equation}
\label{eq:linear-gaussian-law-v04}
X_t\sim\mathcal N\bigl(m(t),P(t)\bigr).
\end{equation}
Для негауссовского $X_0$ среднее и ковариация по-прежнему вычисляются по тем же формулам, но уже не определяют распределение полностью.

\subsection{Постоянные коэффициенты и точный переход по времени}

Для линейного стационарного по коэффициентам SDE
\begin{equation}
\label{eq:lti-sde-v04}
\dd X_t
=
(FX_t+u)\,\dd t
+
G\,\dd W_t
\end{equation}
матрица перехода равна
\begin{equation}
\label{eq:lti-transition-v04}
\Phi(t,s)=\exp\bigl(F(t-s)\bigr).
\end{equation}
Пусть $\Delta t=t_{k+1}-t_k$. Тогда значения процесса на временной сетке имеют точный переходный закон
\begin{equation}
\label{eq:exact-linear-discretization-v04}
X_{t_{k+1}}
=
A_{\Delta t}X_{t_k}+b_{\Delta t}+\xi_k,
\qquad
\xi_k\sim\mathcal N(0,\Sigma_{\Delta t}),
\end{equation}
где
\begin{align}
\label{eq:exact-linear-A-v04}
A_{\Delta t}
&=
\exp(F\Delta t),
\\
\label{eq:exact-linear-b-v04}
b_{\Delta t}
&=
\int_0^{\Delta t}\exp(F(\Delta t-\tau))u\,\dd\tau,
\\
\label{eq:exact-linear-sigma-v04}
\Sigma_{\Delta t}
&=
\int_0^{\Delta t}
\exp(F(\Delta t-\tau))GG^{\trans}
\exp(F(\Delta t-\tau))^{\trans}\,\dd\tau.
\end{align}
Приращения $\xi_k$ независимы на непересекающихся шагах. Равенство~\eqref{eq:exact-linear-discretization-v04} не является методом Эйлера: оно точно воспроизводит конечномерные распределения линейного SDE в узлах сетки. В V05 эта формула станет эталоном для проверки Эйлера--Маруямы.

\begin{computationalnote}[проверка реализации линейного SDE]
Для выбранного шага сначала вычисляют $A_{\Delta t}$ и $\Sigma_{\Delta t}$. Затем проверяют симметрию и неотрицательную определённость $\Sigma_{\Delta t}$, генерируют много переходов из одного состояния и сравнивают выборочные среднее и ковариацию с $A_{\Delta t}x+b_{\Delta t}$ и $\Sigma_{\Delta t}$. Такая проверка отделяет ошибку генератора случайных чисел от ошибки численного интегратора.
\end{computationalnote}

\subsection{Процесс Орнштейна--Уленбека}

Скалярный процесс Орнштейна--Уленбека задаётся уравнением
\begin{equation}
\label{eq:ou-sde-v04}
\dd X_t
=
-\lambda(X_t-\mu)\,\dd t
+
\sigma\,\dd W_t,
\qquad
\lambda>0.
\end{equation}
Параметр $\mu$ задаёт опорный уровень, $\lambda$ --- скорость возврата к нему, а $\sigma$ --- интенсивность случайных возмущений. После сдвига $Y_t=X_t-\mu$ получаем линейное SDE $\dd Y_t=-\lambda Y_t\,\dd t+\sigma\,\dd W_t$.

Точное решение равно
\begin{equation}
\label{eq:ou-solution-v04}
X_t
=
\mu
+
\exp\bigl(-\lambda(t-t_0)\bigr)(X_{t_0}-\mu)
+
\sigma\int_{t_0}^{t}
\exp\bigl(-\lambda(t-\tau)\bigr)\,\dd W_\tau.
\end{equation}
При детерминированном $X_{t_0}=x_0$ условные среднее и дисперсия имеют вид
\begin{align}
\label{eq:ou-mean-v04}
\mathbb E[X_t\mid X_{t_0}=x_0]
&=
\mu+
e^{ -\lambda(t-t_0)}(x_0-\mu),
\\
\label{eq:ou-variance-v04}
\operatorname{Var}(X_t\mid X_{t_0}=x_0)
&=
\frac{\sigma^2}{2\lambda}
\left(1-
e^{-2\lambda(t-t_0)}\right).
\end{align}
Таким образом, переходная плотность нормальна. Среднее забывает начальное отклонение с характерным временем $1/\lambda$, а дисперсия растёт неограниченно лишь при $\lambda=0$. При $\lambda>0$ она стремится к конечному пределу.

Стационарное распределение процесса равно
\begin{equation}
\label{eq:ou-stationary-law-v04}
X_{\infty}
\sim
\mathcal N\!\left(
\mu,
\frac{\sigma^2}{2\lambda}
\right).
\end{equation}
В стационарном режиме ковариация двух моментов зависит только от временного расстояния:
\begin{equation}
\label{eq:ou-autocovariance-v04}
\operatorname{Cov}(X_t,X_{t+\tau})
=
\frac{\sigma^2}{2\lambda}\exp\bigl(-\lambda|\tau|\bigr).
\end{equation}
Следовательно, $\lambda$ управляет не только возвратом среднего, но и длиной памяти процесса.

\subsection{Законченный расчёт: отклонение толщины от рабочего уровня}

Пусть $X_t$ измеряется в миллиметрах и описывает локальную толщину после воздействия случайных производственных возмущений. Возврат к номиналу моделируется параметрами
\begin{equation}
\label{eq:ou-worked-data-v04}
\mu=1.2,
\qquad
\lambda=0.8,
\qquad
\sigma=0.12,
\qquad
X_0=1.5.
\end{equation}
Найдём распределение через $t=2$. Сначала вычисляем множитель памяти
\[
\exp(-\lambda t)
=
\exp(-1.6)
\approx0.2019.
\]
По формуле~\eqref{eq:ou-mean-v04}
\begin{equation}
\label{eq:ou-worked-mean-v04}
m(2)
=
1.2+0.2019(1.5-1.2)
\approx1.2606.
\end{equation}
Дисперсия равна
\begin{equation}
\label{eq:ou-worked-variance-v04}
P(2)
=
\frac{0.12^2}{2\cdot0.8}
\left(1-\exp(-3.2)\right)
\approx0.00863,
\end{equation}
поэтому стандартное отклонение составляет
\[
\sqrt{P(2)}\approx0.0929.
\]
Итак,
\begin{equation}
\label{eq:ou-worked-law-v04}
X_2\sim\mathcal N(1.2606,0.00863).
\end{equation}

Проверим вероятность попасть в технологический интервал $[1.1,1.4]$. После стандартизации получаем
\[
z_{\mathrm{low}}
=
\frac{1.1-1.2606}{0.0929}
\approx-1.73,
\qquad
z_{\mathrm{up}}
=
\frac{1.4-1.2606}{0.0929}
\approx1.50.
\]
Если $\Phi$ --- функция распределения стандартной нормальной величины, то
\begin{equation}
\label{eq:ou-worked-probability-v04}
\mathbb P(1.1\leq X_2\leq1.4)
=
\Phi(1.50)-\Phi(-1.73)
\approx0.9332-0.0418
\approx0.891.
\end{equation}
Вероятность около $89\%$ относится к одной координате и принятой линейно-гауссовской модели. Она не является вероятностью годности всей конструкции без совместной модели остальных параметров и ограничений.

\subsection{Многомерный стационарный режим и ограничения}

Для SDE $\dd X_t=FX_t\,\dd t+G\,\dd W_t$ стационарный гауссовский режим существует, если матрица $F$ устойчива: действительные части всех её собственных значений отрицательны. Тогда среднее стремится к нулю, а стационарная ковариация $P_{\infty}$ решает алгебраическое уравнение Ляпунова
\begin{equation}
\label{eq:continuous-lyapunov-v04}
FP_{\infty}
+
P_{\infty}F^{\trans}
+
GG^{\trans}
=0.
\end{equation}
Это уравнение связывает динамическое подавление и постоянное поступление шума. Если некоторая мода затухает медленно, ковариация вдоль связанного с ней направления обычно велика. Однако собственные векторы $F$ и главные оси $P_{\infty}$ в общем случае не совпадают, особенно для ненормальных матриц.

Для $\tau\geq0$ стационарная временная ковариация задаётся формулами
\begin{equation}
\label{eq:lti-stationary-cross-covariance-v04}
\operatorname{Cov}(X_{t+\tau},X_t)
=
\exp(F\tau)P_{\infty},
\qquad
\operatorname{Cov}(X_t,X_{t+\tau})
=
P_{\infty}\exp(F^{\trans}\tau).
\end{equation}
Они показывают, какие совместные возмущения сохраняются во времени и как связаны разные координаты состояния.

\begin{application}[локальная модель неопределённости конструкции]
Около номинального состояния $x_{\ast}$ нелинейный дрейф можно линеаризовать: $f(x)\approx f(x_{\ast})+F(x-x_{\ast})$, где $F$ --- якобиан. Матрицу $G$ выбирают по физическим источникам неопределённости, а не только для получения удобной диагональной ковариации. После этого уравнение~\eqref{eq:continuous-lyapunov-v04} оценивает типичный совместный разброс геометрии, материала и параметров решётки около рабочего режима.
\end{application}

Линейная модель применима лишь в области, где остаток линеаризации мал. Гауссовское распределение имеет неограниченные хвосты и поэтому само по себе не сохраняет положительность толщины, интервал пористости и геометрическую допустимость. Стационарность невозможна при неустойчивом $F$ без дополнительного управления или нелинейного насыщения. Наконец, плохо масштабированные координаты делают ковариацию и уравнение Ляпунова численно плохо обусловленными. Перед интерпретацией собственных значений $P$ параметры состояния нужно привести к согласованным безразмерным масштабам.

\section{Численное моделирование и метод Эйлера--Маруямы}
\label{sec:numerical-simulation-euler-maruyama}

Для нелинейного SDE обычно невозможно выписать переходную плотность или траекторию в замкнутом виде. Тогда непрерывное время заменяют сеткой, броуновские приращения генерируют явно, а интегральное уравнение аппроксимируют рекурсией. В отличие от обычного ODE, численная схема должна воспроизводить не одну гладкую кривую, а статистику случайных траекторий.

\subsection{Временная сетка и броуновские приращения}

Рассмотрим SDE
\begin{equation}
\label{eq:sde-for-em-v05}
\dd X_t
=
f(X_t,t)\,\dd t
+
G(X_t,t)\,\dd W_t,
\qquad
X_{t_0}=X_0,
\end{equation}
где $X_t\in\R^d$, $W_t\in\R^r$ и $G(x,t)\in\R^{d\times r}$. Выберем сетку
\begin{equation}
\label{eq:time-grid-v05}
t_0<t_1<\cdots<t_N=T,
\qquad
h_k=t_{k+1}-t_k.
\end{equation}
На каждом шаге возникает приращение
\begin{equation}
\label{eq:brownian-increment-grid-v05}
\Delta W_k
=
W_{t_{k+1}}-W_{t_k}
\sim
\mathcal N(0,h_k I_r).
\end{equation}
Приращения для разных $k$ независимы. На практике удобно генерировать
\begin{equation}
\label{eq:brownian-increment-normal-v05}
\Delta W_k=\sqrt{h_k}\,Z_k,
\qquad
Z_k\sim\mathcal N(0,I_r).
\end{equation}
Множитель $\sqrt{h_k}$ обязателен. Замена его на $h_k$ уменьшает дисперсию шага с порядка $h_k$ до порядка $h_k^2$ и меняет моделируемый процесс.

Интегральная форма на одном шаге равна
\begin{equation}
\label{eq:sde-one-step-integral-v05}
X_{t_{k+1}}
=
X_{t_k}
+
\int_{t_k}^{t_{k+1}} f(X_s,s)\,\dd s
+
\int_{t_k}^{t_{k+1}} G(X_s,s)\,\dd W_s.
\end{equation}
Численный метод задаёт правило замены двух интегралов вычислимыми величинами. Поэтому корректность определяется одновременно аппроксимацией дрейфа, диффузии и распределения случайных приращений.

\begin{computationalnote}[воспроизводимость случайной траектории]
Начальное состояние, временная сетка, генератор псевдослучайных чисел и его seed являются частью вычислительного эксперимента. Для сравнения двух схем по одной траектории нужно использовать одну и ту же реализацию броуновского движения, а не два независимых набора шума.
\end{computationalnote}

\subsection{Схема Эйлера--Маруямы}

Если в обоих интегралах из~\eqref{eq:sde-one-step-integral-v05} заморозить коэффициенты в левой точке, получаем
\begin{equation}
\label{eq:euler-maruyama-general-v05}
\widehat X_{k+1}
=
\widehat X_k
+
f(\widehat X_k,t_k)h_k
+
G(\widehat X_k,t_k)\Delta W_k.
\end{equation}
Это метод Эйлера--Маруямы. Для равномерной сетки $h_k=h$ его вычислимая запись имеет вид
\begin{equation}
\label{eq:euler-maruyama-standard-normal-v05}
\widehat X_{k+1}
=
\widehat X_k
+
f(\widehat X_k,t_k)h
+
G(\widehat X_k,t_k)\sqrt h\,Z_k.
\end{equation}
Один шаг требует одного вычисления $f$, одного вычисления $G$ и одного стандартного нормального вектора размерности $r$.

Условно на $\widehat X_k=x$ следующий узел имеет гауссовское распределение
\begin{equation}
\label{eq:euler-maruyama-transition-v05}
\widehat X_{k+1}\mid \widehat X_k=x
\sim
\mathcal N\!\left(
 x+f(x,t_k)h_k,
 G(x,t_k)G(x,t_k)^{\trans}h_k
\right).
\end{equation}
Формула показывает, что схема локально линеаризует не только траекторию, но и переходную плотность. Матрица $G$ не обязана быть квадратной: число независимых источников шума может отличаться от размерности состояния.

Для аддитивного SDE с постоянными $a$ и $G$
\[
\dd X_t=a\,\dd t+G\,\dd W_t
\]
шаг~\eqref{eq:euler-maruyama-general-v05} совпадает с точным переходом. Для линейного SDE с матричным дрейфом он уже приближённый; точный переход задан формулами~\eqref{eq:exact-linear-discretization-v04}--\eqref{eq:exact-linear-sigma-v04}.

\begin{application}[зашумление параметров конструкции]
Пусть $X_k$ содержит нормированные координаты геометрии, материала, толщины и решётки. Тогда $f(X_k,t_k)h$ задаёт направленное изменение за шаг, а $G(X_k,t_k)\sqrt h Z_k$ --- случайное изменение с физически заданными корреляциями. После обновления нужно отдельно проверить допустимость декодированной конструкции; сама схема не гарантирует положительность толщины, отсутствие самопересечений или выполнение ограничений симулятора.
\end{application}

\subsection{Сильная и слабая точность}

Для случайного метода недостаточно одного понятия ошибки. Пусть точное и численное решения построены на одном броуновском движении и сравниваются в момент $T$.

\begin{definition}[сильная сходимость]
Метод имеет сильный порядок $\gamma$, если при достаточно малом максимальном шаге $h=\max_k h_k$ выполняется оценка
\begin{equation}
\label{eq:strong-convergence-v05}
\mathbb E\!
\left[
\lVert X_T-\widehat X_N\rVert
\right]
\leq C h^{\gamma}.
\end{equation}
\end{definition}
Сильная ошибка измеряет близость отдельных траекторий при одинаковом шуме. Она важна, когда результат зависит от пути: момента выхода за границу, максимальной деформации, накопленного повреждения или сопряжения SDE с физическим симулятором.

\begin{definition}[слабая сходимость]
Метод имеет слабый порядок $\alpha$, если для достаточно гладкой тестовой функции $\varphi$ выполняется
\begin{equation}
\label{eq:weak-convergence-v05}
\left|
\mathbb E[\varphi(X_T)]
-
\mathbb E[\varphi(\widehat X_N)]
\right|
\leq C_{\varphi} h^{\alpha}.
\end{equation}
\end{definition}
Слабая ошибка относится к распределению и его функционалам. Она подходит для оценки среднего целевого функционала, вероятности события после гладкого приближения индикатора и статистики ансамбля.

При стандартных условиях регулярности метод Эйлера--Маруямы имеет сильный порядок $1/2$ и слабый порядок $1$. При аддитивном шуме, когда $G$ не зависит от состояния, сильный порядок повышается до $1$. Эти порядки являются асимптотическими: на грубой сетке ошибка может не следовать степенному закону.

Сильную и слабую задачи нельзя смешивать. Две схемы могут давать близкие средние, но заметно разные траектории. Обратная ситуация также возможна для негладкого функционала, например вероятности редкого выхода за ограничение.

\begin{computationalnote}[согласование сеток]
При сравнении шагов $h$ и $h/2$ одну реализацию броуновского движения строят на мелкой сетке. Крупное приращение получают суммой двух мелких:
\begin{equation}
\label{eq:brownian-coupling-v05}
\Delta W_k^{(h)}
=
\Delta W_{2k}^{(h/2)}
+
\Delta W_{2k+1}^{(h/2)}.
\end{equation}
Тогда разность траекторий отражает ошибку дискретизации, а не замену шума.
\end{computationalnote}

\subsection{Законченный расчёт: два шага для толщины}

Рассмотрим процесс возврата толщины к номиналу
\begin{equation}
\label{eq:ou-em-example-sde-v05}
\dd X_t
=
-2(X_t-1)\,\dd t
+
0.1\,\dd W_t,
\qquad
X_0=1.2.
\end{equation}
Возьмём шаг $h=0.25$ и два стандартных нормальных значения
\begin{equation}
\label{eq:ou-em-example-noise-v05}
Z_0=0.4,
\qquad
Z_1=-0.7.
\end{equation}
По~\eqref{eq:brownian-increment-normal-v05}
\[
\Delta W_0=\sqrt{0.25}\cdot0.4=0.2,
\qquad
\Delta W_1=\sqrt{0.25}\cdot(-0.7)=-0.35.
\]
Первый шаг Эйлера--Маруямы равен
\begin{align}
\widehat X_1
&=
1.2-2(1.2-1)\cdot0.25+0.1\cdot0.2
\notag\\
&=
1.2-0.1+0.02
=
1.12.
\label{eq:ou-em-example-step1-v05}
\end{align}
Второй шаг даёт
\begin{align}
\widehat X_2
&=
1.12-2(1.12-1)\cdot0.25+0.1\cdot(-0.35)
\notag\\
&=
1.12-0.06-0.035
=
1.025.
\label{eq:ou-em-example-step2-v05}
\end{align}
Значение $1.025$ является одной случайной реализацией, а не средним и не наиболее вероятным результатом.

Теперь сравним распределения в момент $T=0.5$. Точная формула~\eqref{eq:ou-mean-v04} даёт
\begin{equation}
\label{eq:ou-em-exact-mean-v05}
\mathbb E[X_{0.5}]
=
1+0.2e^{-1}
\approx1.0736,
\end{equation}
а по~\eqref{eq:ou-variance-v04}
\begin{equation}
\label{eq:ou-em-exact-variance-v05}
\operatorname{Var}(X_{0.5})
=
\frac{0.1^2}{4}\left(1-e^{-2}\right)
\approx0.00216.
\end{equation}

Для схемы обозначим $Y_k=\widehat X_k-1$. Тогда
\begin{equation}
\label{eq:ou-em-ar1-v05}
Y_{k+1}
=
(1-2h)Y_k+0.1\sqrt h\,Z_k
=
0.5Y_k+0.05Z_k.
\end{equation}
После двух шагов
\begin{equation}
\label{eq:ou-em-two-step-mean-v05}
\mathbb E[\widehat X_2]
=
1+0.5^2(1.2-1)
=
1.05,
\end{equation}
и
\begin{equation}
\label{eq:ou-em-two-step-variance-v05}
\operatorname{Var}(\widehat X_2)
=
0.1^2 h\left(1+0.5^2\right)
=
0.003125.
\end{equation}
Грубая сетка занизила среднее на $0.0236$ и завысила дисперсию примерно на $45\%$. Совпадение одной траектории с допустимым интервалом не устраняет этот систематический сдвиг распределения.

\subsection{Устойчивость, шаг и смещение стационарной дисперсии}

Для общего процесса Орнштейна--Уленбека
\[
\dd X_t=-\lambda(X_t-\mu)\,\dd t+\sigma\,\dd W_t
\]
схема принимает вид
\begin{equation}
\label{eq:ou-em-recursion-v05}
\widehat X_{k+1}-\mu
=
(1-\lambda h)(\widehat X_k-\mu)
+
\sigma\sqrt h\,Z_k.
\end{equation}
Детерминированная часть затухает только при
\begin{equation}
\label{eq:ou-em-stability-v05}
|1-\lambda h|<1,
\qquad
0<h<\frac{2}{\lambda}.
\end{equation}
Это условие устойчивости, а не точности. Шаг, близкий к $2/\lambda$, может формально давать ограниченную последовательность, но искажать среднее, ковариацию и временную корреляцию.

Если~\eqref{eq:ou-em-stability-v05} выполнено, стационарная дисперсия дискретной схемы равна
\begin{equation}
\label{eq:ou-em-stationary-variance-v05}
P_{\infty}^{\mathrm{EM}}
=
\frac{\sigma^2h}
{1-(1-\lambda h)^2}
=
\frac{\sigma^2}{2\lambda-\lambda^2h}.
\end{equation}
Точная дисперсия процесса равна $P_{\infty}=\sigma^2/(2\lambda)$. Поэтому явная схема систематически завышает её при любом $h>0$.

Для параметров расчёта выше
\[
P_{\infty}=\frac{0.01}{4}=0.0025,
\qquad
P_{\infty}^{\mathrm{EM}}
=
\frac{0.01}{4-4\cdot0.25}
\approx0.00333.
\]
Относительное завышение составляет около $33\%$. Если линейный переход известен точно, использование~\eqref{eq:exact-linear-discretization-v04} устраняет эту дискретизационную ошибку при тех же узлах времени.

Жёсткий дрейф наследует ограничения явного Эйлера из главы~\ref{chap:dynamics-differentiable-models}. Малый шум не делает крупный шаг безопасным. При мультипликативном шуме дополнительно возможны редкие большие выбросы, отрицательные значения положительной переменной и рост моментов, незаметный по нескольким траекториям.

\subsection{Проверка симулятора и практические ограничения}

Проверку реализации начинают с модели, для которой известен точный переход. Для линейного SDE из V04 из одного состояния генерируют большой ансамбль шагов и сравнивают выборочные среднее и ковариацию с аналитическими. Затем повторяют расчёт на сетках $h$, $h/2$ и $h/4$ с согласованными броуновскими приращениями.

Для $M$ траекторий конечные среднее и ковариация оцениваются как
\begin{align}
\label{eq:monte-carlo-mean-v05}
\widehat m_T
&=
\frac1M\sum_{j=1}^{M}\widehat X_N^{(j)},
\\
\label{eq:monte-carlo-covariance-v05}
\widehat P_T
&=
\frac1{M-1}
\sum_{j=1}^{M}
\bigl(\widehat X_N^{(j)}-\widehat m_T\bigr)
\bigl(\widehat X_N^{(j)}-\widehat m_T\bigr)^{\trans}.
\end{align}
Общая ошибка содержит две части: смещение по шагу и статистическую ошибку конечного ансамбля. Уменьшение $h$ не компенсирует слишком малое $M$, а увеличение $M$ не исправляет грубую дискретизацию.

Простое отсечение недопустимых координат после каждого шага меняет переходный закон и обычно создаёт смещение около границы. Проекция на многообразие также задаёт уже другую дискретную динамику. Если ограничения существенны, предпочтительны подходящая параметризация, отражающие или поглощающие границы либо специально построенное геометрическое SDE.

В задачах генеративного проектирования одна траектория может требовать многократного вызова дорогого физического симулятора. Тогда $h$ выбирают не только по стоимости, но и по проверенной ошибке интересующего функционала: массы, прочности, потери давления или нарушения ограничения. Координаты состояния нормируют до интегрирования; иначе единый шаг может быть приемлем для одной группы параметров и разрушителен для другой.

При диффузионном обучении прямой процесс создаёт зашумлённые состояния, а обратное SDE генерирует кандидатов. Если переходное распределение прямого линейного SDE известно аналитически, обучающие пары выгоднее получать прямой выборкой из него, не накапливая ошибку Эйлера--Маруямы. Для общего обратного SDE численный решатель неизбежен, и рассмотренные различия между траекторной и распределительной точностью переходят в V06.

\section{Обратное SDE, score-функция и диффузионные генеративные модели}
\label{sec:reverse-sde-score-diffusion-models}

Предыдущие разделы описывали прямую задачу: заданное SDE порождает случайные траектории и переносит их распределение. В генеративной постановке исходным является распределение допустимых конструкций $p_{\mathrm{data}}$. Нужно построить простой прямой процесс, который постепенно уничтожает его структуру, а затем научиться обращать эту эволюцию от известного шума к новым образцам.

\subsection{Прямое зашумление: от данных к простому распределению}

Пусть $X_0\sim p_{\mathrm{data}}$ --- нормированное непрерывное представление конструкции. Прямой процесс Song и соавторов задаётся SDE
\begin{equation}\label{eq:forward-generative-sde-v06}
\dd X_t = f(X_t,t)\,\dd t + g(t)\,\dd W_t, \qquad 0\leq t\leq T.
\end{equation}
Коэффициенты выбирают так, чтобы распределение $p_T$ было простым и доступным для прямой генерации, обычно гауссовским. Нейронная сеть в~\eqref{eq:forward-generative-sde-v06} не участвует: закон зашумления фиксируется заранее.

При аффинном дрейфе переход $p_{0t}(x_t\mid x_0)$ является гауссовским. Удобная форма имеет вид
\begin{equation}\label{eq:gaussian-perturbation-kernel-v06}
X_t = \alpha(t)X_0 + \sigma(t)Z, \qquad Z\sim\mathcal N(0,I_d).
\end{equation}
Она позволяет получать обучающую пару $(X_0,X_t)$ сразу, без последовательного применения метода Эйлера--Маруямы. Функция $\alpha$ задаёт сохранённую долю исходного сигнала, а $\sigma$ --- масштаб добавленного шума.

В variance-exploding процессе среднее условного перехода сохраняется, а дисперсия растёт. В variance-preserving процессе дрейф одновременно уменьшает исходный сигнал, поэтому дисперсия остаётся контролируемой. Выбор между ними является частью модели: он влияет на масштабы входов, сложность оценки score и численную устойчивость обратной динамики.

Для генеративного проектирования вектор $X_0$ не обязан состоять из пикселей. Он может кодировать координаты фиксированной сетки, поле плотности, толщины элементов, параметры материала или латентное представление геометрии. Однако все координаты должны образовывать непрерывное евклидово пространство фиксированной размерности. Переменная связность сетки, категориальные материалы и дискретная топология требуют отдельной кодировки; добавление гауссовского шума к их исходным индексам математического смысла не имеет.

\subsection{Score-функция и её обучение}

\begin{definition}
Score-функцией плотности $p_t$ называется векторное поле
\begin{equation}\label{eq:score-definition-v06}
s_t(x) = \grad_x\log p_t(x).
\end{equation}
\end{definition}
В точке $x$ score указывает направление наиболее быстрого локального роста логарифма плотности. Он не равен вероятности и не задаёт расстояние до ближайшего образца. Нормирующая константа плотности исчезает после дифференцирования, поэтому поле можно оценивать без явного вычисления $p_t(x)$.

Модель $s_\theta(x,t)$ обучают приближать~\eqref{eq:score-definition-v06} для всех уровней шума. В denoising score matching используется функционал
\begin{equation}\label{eq:denoising-score-objective-v06}
\mathcal L(\theta) = \mathbb E \left[ \lambda(t) \left\lVert s_\theta(X_t,t) - \grad_{x_t}\log p_{0t}(X_t\mid X_0) \right\rVert^2 \right],
\end{equation}
где усреднение проводится по времени, данным $X_0$ и зашумлённому состоянию $X_t$. При достаточной ёмкости модели оптимум~\eqref{eq:denoising-score-objective-v06} совпадает с маргинальным score $\grad_x\log p_t(x)$ почти всюду.

Для ядра~\eqref{eq:gaussian-perturbation-kernel-v06}
\begin{equation}\label{eq:gaussian-conditional-score-v06}
\grad_{x_t}\log p_{0t}(x_t\mid x_0) = - \frac{x_t-\alpha(t)x_0}{\sigma^2(t)}.
\end{equation}
Следовательно, обучающая цель вычисляется непосредственно по чистому образцу, времени и сгенерированному шуму. Сеть восстанавливает не сам $X_0$, а поле направлений для распределения на каждом $t$.

Score зависит от выбора координат. Если нормированная переменная задаётся как $y=S^{-1}(x-\mu)$, то правило цепочки из главы~\ref{chap:multivariable-calculus} даёт
\begin{equation}\label{eq:score-rescaling-v06}
s_x(x) = S^{-\trans}s_y(y).
\end{equation}
Поэтому ненормированные миллиметры, паскали и безразмерная пористость создают искусственное преобладание отдельных координат. Матрицу масштабирования нужно сохранять вместе с моделью и учитывать при физическом управлении генерацией.

\subsection{Обратное SDE и численная генерация}

Для скалярного коэффициента диффузии, не зависящего от состояния, обратный процесс имеет вид
\begin{equation}\label{eq:reverse-time-sde-v06}
\dd X_t = \left[ f(X_t,t) - g^2(t)s_t(X_t) \right]\dd t + g(t)\,\dd\overline W_t,
\end{equation}
где время идёт от $T$ к нулю, а $\overline W_t$ является броуновским движением в обратном времени. В вычислениях неизвестный $s_t$ заменяется сетью $s_\theta$.

Пусть $t_{k-1}=t_k-h$ и $h>0$. Один обратный шаг Эйлера--Маруямы записывается как
\begin{equation}\label{eq:reverse-em-step-v06}
X_{k-1} = X_k - \left[ f(X_k,t_k) - g^2(t_k)s_\theta(X_k,t_k) \right]h + g(t_k)\sqrt h\,Z_k, \qquad Z_k\sim\mathcal N(0,I_d).
\end{equation}
Знак перед квадратной скобкой появляется потому, что шаг времени отрицателен. Начальное значение для генерации берётся из простого распределения $X_T\sim p_T$.

Score изменяет дрейф так, чтобы компенсировать рассеивание прямого процесса. Он не отменяет случайность: разные $Z_k$ при одном $X_T$ дают разные траектории. Ошибка генератора состоит как минимум из ошибки сети, ошибки временной дискретизации и несоответствия конечного $p_T$ выбранному простому опорному распределению.

В predictor--corrector схеме шаг~\eqref{eq:reverse-em-step-v06} служит предиктором. Затем при фиксированном времени выполняют один или несколько шагов Ланжевена
\begin{equation}\label{eq:langevin-corrector-v06}
X \leftarrow X + \varepsilon s_\theta(X,t) + \sqrt{2\varepsilon}\,Z.
\end{equation}
Предиктор переносит состояние между соседними распределениями, а корректор уменьшает ошибку внутри текущего $p_t$. Дополнительные шаги повышают стоимость, поэтому их число и $\varepsilon$ проверяют по качеству образцов и числу вызовов score-сети, а не выбирают только по визуальному результату нескольких траекторий.

\subsection{Probability-flow ODE и связь с дифференцируемой динамикой}

Тому же семейству маргинальных плотностей $p_t$ соответствует детерминированное уравнение
\begin{equation}\label{eq:probability-flow-ode-v06}
\frac{\dd X_t}{\dd t} = f(X_t,t) - \frac12 g^2(t)s_t(X_t).
\end{equation}
Оно называется probability-flow ODE. Его отдельные траектории не совпадают с траекториями SDE, но при точном score распределение $X_t$ в каждый момент совпадает с распределением прямого или обратного стохастического процесса.

После замены $s_t$ на $s_\theta$ правая часть~\eqref{eq:probability-flow-ode-v06} становится обучаемым векторным полем, изученным в главе~\ref{chap:dynamics-differentiable-models}. Можно использовать адаптивный ODE-решатель, контролировать локальную погрешность и дифференцировать результат по параметрам. Плата за детерминированность состоит в том, что каждая траектория однозначно задаётся начальным шумом и допуском решателя.

Обозначим
\begin{equation}\label{eq:probability-flow-vector-field-v06}
v_\theta(x,t) = f(x,t) - \frac12 g^2(t)s_\theta(x,t).
\end{equation}
Вдоль probability-flow ODE изменение логарифма плотности выражается через дивергенцию:
\begin{equation}\label{eq:probability-flow-log-density-v06}
\frac{\dd}{\dd t}\log p_t(X_t) = -\div_x v_\theta(X_t,t).
\end{equation}
Это связывает score-модель с непрерывными нормализующими потоками и позволяет оценивать правдоподобие. В большой размерности точное вычисление дивергенции дорого; обычно требуется стохастическая оценка следа, которая добавляет вычислительную дисперсию.

Равенство маргиналов является теоретическим свойством точного поля. При неточном $s_\theta$, грубом допуске решателя или ошибочной нормировке SDE и ODE могут давать заметно разные распределения. Поэтому совпадение нескольких изображений или конструкций не заменяет статистического сравнения ансамблей.

\subsection{Законченный расчёт: score и один обратный шаг}

Рассмотрим одномерный прямой процесс
\begin{equation}\label{eq:ve-calculation-forward-v06}
\dd X_t=\dd W_t, \qquad X_0\sim\mathcal N(1,0.25).
\end{equation}
По формулам V02 и V04
\begin{equation}\label{eq:ve-calculation-marginal-v06}
X_t\sim\mathcal N(1,0.25+t).
\end{equation}
В момент $t=0.5$ дисперсия равна $0.75$, поэтому маргинальный score имеет вид
\begin{equation}\label{eq:ve-calculation-score-v06}
s_{0.5}(x) = -\frac{x-1}{0.75}.
\end{equation}
Для зашумлённого состояния $x=0.2$
\begin{equation}\label{eq:ve-calculation-score-value-v06}
s_{0.5}(0.2) = -\frac{0.2-1}{0.75} \approx1.0667.
\end{equation}
Положительный score указывает в сторону центра распределения $x=1$.

Здесь $f=0$ и $g=1$, поэтому обратный дрейф в~\eqref{eq:reverse-time-sde-v06} равен
\begin{equation}\label{eq:ve-calculation-reverse-drift-v06}
b_{\mathrm{rev}}(0.2,0.5) = -s_{0.5}(0.2) \approx-1.0667.
\end{equation}
Возьмём обратный шаг $h=0.05$ и сначала отделим детерминированную часть, положив $Z=0$. По~\eqref{eq:reverse-em-step-v06}
\begin{equation}\label{eq:ve-calculation-reverse-step-v06}
X_{0.45} \approx 0.2-(-1.0667)\cdot0.05 \approx 0.2533.
\end{equation}
Хотя коэффициент обратного дрейфа отрицателен, отрицательный шаг времени перемещает точку вправо, к области большей плотности. При полном шаге добавляется $\sqrt{0.05}Z$, поэтому одна реализация может двигаться и в противоположную сторону; правильный эффект наблюдается на распределении траекторий.

Для отдельной обучающей пары условный переход равен $X_t\mid X_0=x_0\sim\mathcal N(x_0,t)$, и цель~\eqref{eq:gaussian-conditional-score-v06} принимает вид $-(x_t-x_0)/t$. Эта условная величина не обязана совпадать с~\eqref{eq:ve-calculation-score-value-v06} для одного $x_0$: маргинальный score возникает после усреднения по чистым образцам, совместимым с наблюдаемым $x_t$.

\subsection{Генерация конструкции, управление и ограничения}

Практический конвейер начинается с кодировки $q\mapsto X_0$, где $q$ описывает геометрию и материал, а $X_0$ является нормированным вектором фиксированной размерности. Прямое SDE создаёт пары разных уровней шума, сеть учится оценивать $s_t$, после чего обратное SDE или probability-flow ODE переводит случайное опорное распределение в кандидата $\widehat X_0$. Декодер возвращает физическую конструкцию, которую независимо проверяют симулятором и технологическими ограничениями.

Условная генерация использует дополнительное наблюдение или требование $y$. В формуле Song и соавторов к безусловному score добавляется
\begin{equation}\label{eq:conditional-score-v06}
\grad_x\log p_t(x\mid y) = \grad_x\log p_t(x) + \grad_x\log p_t(y\mid x).
\end{equation}
Второй член направляет обратную динамику к состояниям, согласованным с условием. Для генеративного проектирования $y$ может кодировать класс конструкции, заданную область, нагрузочный сценарий или измеренные характеристики.

Если физический симулятор дифференцируем, прикладную модель правдоподобия можно задать, например, как
\begin{equation}\label{eq:physics-guidance-likelihood-v06}
\log p(y\mid x) = -\kappa J_{\mathrm{phys}}(x) - \sum_j \rho_j\,\phi\bigl(c_j(x)\bigr),
\end{equation}
где $J_{\mathrm{phys}}$ --- целевой функционал, $c_j(x)\leq0$ --- ограничения, а $\phi$ --- гладкий штраф. Тогда градиенты из глав~\ref{chap:multivariable-calculus} и~\ref{chap:dynamics-differentiable-models} входят в условный score. Это не универсальная формула из обучения данных, а явный выбор вероятностной модели требований; коэффициенты $\kappa$ и $\rho_j$ определяют силу управления и должны иметь согласованный масштаб.

Слишком сильное управление может вывести траекторию в область, где score-сеть не обучалась, и создать формально оптимальный, но геометрически бессмысленный объект. Жёсткие ограничения надёжнее обеспечивать параметризацией или декодером, а не только штрафом. Проекция после каждого шага меняет обратную динамику и требует отдельной проверки смещения.

\begin{application}[минимальная проверка генеративной модели конструкции]
Нужно отдельно проверить качество прямого ядра, ошибку score на удержанной выборке, сходимость обратного решателя по шагу, долю допустимых декодированных объектов и распределение физических метрик. Затем кандидаты сравнивают с обучающим набором, чтобы отличить генерацию новых решений от воспроизведения ближайших примеров. Финальная оценка проводится точным или более дорогим симулятором, не участвовавшим в управлении.
\end{application}

Score определён через плотность в евклидовом пространстве. Если данные лежат на тонком многообразии, при $t>0$ шум делает плотность гладкой, но задача становится плохо обусловленной при $t\to0$. Для сеток с меняющейся топологией, дискретных материалов и разрывных ограничений обычная формула требует другой модели состояния. Диффузионная генерация также не гарантирует физическую корректность, разнообразие и превосходство над исходной выборкой: она аппроксимирует распределение, заданное данными и выбранным управлением.

Итоговая прикладная цепочка главы теперь замкнута:
\begin{equation}\label{eq:chapter-five-final-chain-v06}
\begin{aligned} \text{неопределённость} &\longrightarrow \text{SDE} \longrightarrow p_t \longrightarrow \grad\log p_t, \\ \grad\log p_t &\longrightarrow \text{обратная динамика} \longrightarrow \text{кандидат конструкции}. \end{aligned}
\end{equation}
Каждая стрелка требует собственной проверки: постановки пространства состояния, корректности стохастической модели, качества score, численной точности и физической допустимости результата.
